% Vietnamese translation for OI Public Library
% Source-PDF: IOI中国国家候选队论文/2009/徐源盛/对一类动态规划问题的研究.pdf
\documentclass[11pt]{article}
\usepackage{oipl-vn}

\newcommand{\Oh}{\mathcal{O}}

\title{Nghiên cứu về một lớp bài toán quy hoạch động}
\author{Từ Nguyên Thịnh\\Trường Trung học số 1 Trường Sa, Hồ Nam}
\date{2009}

\begin{document}
\maketitle

\origtitle{A study of a class of dynamic programming problems}
\sourcepath{IOI-China-Candidate-Team-Papers/2009/Xu-Yuansheng/future-cost-dp.pdf}

\begin{abstract}
Bài viết khảo sát, thông qua bốn bài toán, một lớp bài toán quy hoạch động khá
đặc biệt: quyết định ở hiện tại ảnh hưởng đến chi phí của ``hành động'' trong
tương lai. Nếu ảnh hưởng của quyết định hiện tại lên tương lai chỉ phụ thuộc
vào chính quyết định hiện tại, ta trực tiếp tính ảnh hưởng đó vào chi phí của
quyết định hiện tại rồi truyền qua trạng thái. Nếu ảnh hưởng lên tương lai còn
phụ thuộc vào tình huống tương lai, ta bổ sung trạng thái để giả định tình
huống tương lai; khi đi đến quyết định tương lai thì dùng thẳng trạng thái đã
giả định đó. Đây là phương pháp giải được trình bày chi tiết trong bài.
\end{abstract}

\noindent\textbf{Từ khóa:} quy hoạch động, tính trước chi phí, giả định quyết
định tương lai.

\section{Mở đầu}

Trong các bài toán quy hoạch động thông thường, chi phí do ``hành động'' ở
trạng thái hiện tại thường được tính đồng thời với trạng thái đó. Chẳng hạn
trong bài toán gộp đá kinh điển, phương trình là
\[
  f[i][j]=\max_k\{f[i][k]+f[k+1][j]\}+w(i,j).
\]
Khi tính $f[i][j]$, ta mới cộng chi phí của hành động gộp toàn bộ đống đá từ
$i$ đến $j$. Điều này rất phù hợp với thói quen tư duy.

Tuy nhiên những năm gần đây thường xuất hiện một lớp bài toán quy hoạch động
trong đó một phần chi phí của ``hành động'' hiện tại phải được tính từ các
quyết định trước đó, rồi thông qua trạng thái mà tác động đến trạng thái hiện
tại. Nguyên nhân của cách tính đặc biệt này là quyết định hiện tại có ảnh
hưởng đến chi phí của ``hành động'' trong tương lai. Việc xây dựng phương trình
cho lớp bài toán này thường khó hơn, cần phân tích kỹ đề bài và tìm ra mâu
thuẫn cốt lõi.

\section{Quyết định hiện tại chỉ ảnh hưởng đến tương lai qua chính nó}

Xét một ví dụ đơn giản: trong hai nam và hai nữ cần chọn một nam một nữ đi làm
nhiệm vụ, biết hiệu suất của từng người và muốn tổng hiệu suất lớn nhất. Nếu
nam $A$ được chọn thì hiệu suất của mọi nữ tăng $7$; nếu nam $B$ được chọn thì
hiệu suất của mọi nữ giảm $7$. Ở giai đoạn thứ nhất, chọn $A$ hay $B$ có ảnh
hưởng khác nhau đến hiệu suất của nữ ở giai đoạn thứ hai. Ta có thể xem ảnh
hưởng này là ``sức hút cá nhân'' của nam: cộng $7$ vào hiệu suất của $A$ và
trừ $7$ khỏi hiệu suất của $B$. Thực chất, ta đã tính chi phí của giai đoạn
thứ hai ngay ở giai đoạn thứ nhất. Với dạng bài như vậy, ta thường gộp ảnh
hưởng lên ``hành động'' tương lai vào chi phí của quyết định hiện tại. Hai ví
dụ sau trình bày chi tiết hơn ý tưởng này.

\subsection{Bài toán 1: Những quả cầu của Sue (SDTSC 2008)}

\subsubsection*{Đề bài}

Sue và Sandy gần đây say mê một trò chơi máy tính. Câu chuyện trong trò chơi
diễn ra trên biển cả đẹp, bí ẩn và đầy kích thích. Sue có một chiếc thuyền nhỏ
nhẹ, nhưng mục tiêu của Sue không phải làm hải tặc, mà là thu thập những quả
trứng màu đang trôi nổi trong không trung. Sue có một vũ khí bí mật: chỉ cần
chèo thuyền đến đúng phía dưới một quả trứng rồi dùng vũ khí đó, cô có thể thu
thập quả trứng ngay lập tức.

Mỗi quả trứng có một giá trị hấp dẫn, và giá trị này giảm dần theo thời gian
rơi trong không trung. Muốn được nhiều điểm hơn, Sue phải cố thu thập quả
trứng khi giá trị hấp dẫn còn cao. Nếu một quả trứng rơi xuống biển, giá trị
hấp dẫn của nó trở thành một số âm, nhưng điều này không làm Sue mất hứng, vì
mỗi quả trứng đều khác nhau và Sue muốn thu thập tất cả.

Sandy thì không lãng mạn như Sue; Sandy muốn đạt nhiều điểm nhất có thể. Để
giải quyết bài toán, trước hết Sandy trừu tượng hóa trò chơi thành mô hình
sau.

Lấy mặt phẳng ngang chứa vị trí ban đầu của Sue làm trục $x$. Ban đầu có $N$
quả trứng trên không. Với quả trứng thứ $i$, vị trí ban đầu là tọa độ nguyên
$(x_i,y_i)$. Sau khi trò chơi bắt đầu, nó rơi đều theo chiều âm của trục $y$
với vận tốc $v_i$ đơn vị khoảng cách trên một đơn vị thời gian. Vị trí ban đầu
của Sue là $(x_0,0)$. Sue có thể di chuyển theo chiều dương hoặc âm của trục
$x$, với vận tốc $1$ đơn vị khoảng cách trên một đơn vị thời gian. Việc dùng
vũ khí bí mật để lấy một quả trứng là tức thời; điểm nhận được bằng một phần
nghìn tung độ hiện tại của quả trứng.

Hãy giúp Sue và Sandy: với điều kiện phải thu thập mọi quả trứng, hãy làm cho
tổng điểm lớn nhất.

\subsubsection*{Dữ liệu vào}

Dòng đầu gồm hai số nguyên $N,x_0$, cách nhau bởi một dấu cách, lần lượt là số
quả trứng và vị trí ban đầu của Sue. Dòng thứ hai gồm $N$ số nguyên $x_i$.
Dòng thứ ba gồm $N$ số nguyên $y_i$. Dòng thứ tư gồm $N$ số nguyên $v_i$, là
vận tốc rơi đều theo chiều âm của trục $y$.

\subsubsection*{Dữ liệu ra}

In một số thực, làm tròn đến ba chữ số sau dấu thập phân: điểm lớn nhất có thể
đạt được khi đã thu thập tất cả quả trứng.

\subsubsection*{Ví dụ}

\begin{verbatim}
ball.in
3 0
-4 -2 2
22 30 26
1 9 8

ball.out
0.000
\end{verbatim}

\subsubsection*{Ràng buộc}

Với $30\%$ dữ liệu, $N\le 20$. Với $100\%$ dữ liệu, $N\le 1000$ và
$-10^4\le x_i,y_i,v_i\le 10^4$.

\subsubsection*{Phân tích}

Trước hết sắp xếp tất cả điểm, bao gồm cả điểm xuất phát, theo hoành độ. Khi
đó bài toán trở thành: xuất phát từ điểm ban đầu, mỗi lần đi sang trái hoặc
sang phải để đến một quả trứng gần nhất chưa lấy và bắn rơi nó. Gọi $t_i$ là
thời điểm lần đầu tiên đi đến quả trứng thứ $i$, đáp án cần tối đa hóa là
\[
  \sum_i (y_i-t_i v_i).
\]

Ta nhận thấy tập các quả trứng đã bắn rơi luôn mở rộng thành một đoạn. Vì thế
rất tự nhiên nghĩ đến $f_1[i][j]$ và $f_2[i][j]$: lần lượt là điểm lớn nhất
khi đã bắn rơi đoạn quả trứng từ $i$ đến $j$, hiện đang dừng ở điểm $i$ hoặc
điểm $j$. Xét $f_1[i][j]$, tức điểm $i$ là quả trứng vừa bắn. Điểm nhận được
khi bắn phụ thuộc vào thời điểm hiện tại, nhưng thời điểm này không thể biểu
diễn trong trạng thái $f_1[i][j]$. Nếu thêm một chiều trạng thái để biểu diễn
thời gian, độ phức tạp sẽ không chấp nhận được.

Mâu thuẫn ở đây là phương án di chuyển trong quá khứ ảnh hưởng đến điểm bắn ở
hiện tại. Hãy xét kỹ hơn cách tính $f_1[i][j]$. Điểm bắn quả trứng $i$ là
$y_i-t v_i$, trong đó $t$ bằng tổng thời gian di chuyển của mọi quyết định
trước đó. Do đó, giống ví dụ chọn nam nữ ở trên, ta có thể tính phần $-t v_i$
ngay trong các bước di chuyển trước. Nói cách khác, mỗi lần di chuyển phải
tính luôn lượng điểm chắc chắn sẽ mất trong tương lai. Chẳng hạn, khi chuyển
từ $f_1[i][j]$ sang $f_1[i-1][j]$, tức đi từ $i$ đến $i-1$, ngoài đoạn quả
trứng $i$ đến $j$ đã bắn rơi, mọi quả trứng khác đều đang rơi; lượng điểm bị
mất này được tính vào chi phí đi từ $i$ đến $i-1$. Vì phần $-t v_i$ đã được
tính ở các quyết định trước, lúc bắn chỉ cần cộng $y_i$.

Để mô tả thuận tiện, đặt
\[
  w[i][j]=\sum_{r=1}^{n} v_r-\sum_{k=i}^{j} v_k,
\]
tức tổng vận tốc rơi của mọi quả trứng nằm ngoài đoạn $i..j$. Khi đó dễ thu
được phương trình
\[
\begin{aligned}
f_1[i][j] = y_i+\max\{&
  f_1[i+1][j]-(x_{i+1}-x_i)w[i+1][j],\\
& f_2[i+1][j]-(x_j-x_i)w[i+1][j]\},
\end{aligned}
\]
\[
\begin{aligned}
f_2[i][j] = y_j+\max\{&
  f_2[i][j-1]-(x_j-x_{j-1})w[i][j-1],\\
& f_1[i][j-1]-(x_j-x_i)w[i][j-1]\}.
\end{aligned}
\]
Đáp án là
\[
  \frac{\max\{f_1[1][n],f_2[1][n]\}}{1000}.
\]
Độ phức tạp thời gian là $\Oh(n^2)$. Đến đây bài toán được giải quyết hoàn
toàn.

Nhìn lại quá trình trên, điểm số khi bắn quả trứng đã được tách thành hai phần
để tính. Việc toàn bộ quả trứng rơi xuống chỉ phụ thuộc vào thời gian di
chuyển hiện tại, nên ta có thể tính trước điểm tương lai vào trạng thái hiện
tại.

\subsection{Bài toán 2: Bài toán sắp xếp thứ hạng (BOI 2007, ngày 1)}

\subsubsection*{Đề bài}

Biết điểm số của các thí sinh, nhiệm vụ là đưa ra thứ hạng của họ theo thứ tự
điểm từ cao xuống thấp. Đáng tiếc, cấu trúc dữ liệu lưu thông tin thí sinh chỉ
hỗ trợ một thao tác: chuyển một thí sinh từ vị trí $i$ đến vị trí $j$. Thao
tác này không làm thay đổi thứ tự tương đối của các thí sinh khác. Nếu
$i>j$, các thí sinh ở vị trí từ $j$ đến $i-1$ đều lùi về sau một vị trí; nếu
$i<j$, các thí sinh ở vị trí từ $i+1$ đến $j$ đều tiến lên trước một vị trí.
Chi phí của thao tác chuyển là $i+j$, trong đó vị trí được đánh số từ $1$.

Hãy xác định một dãy thao tác để sắp xếp thí sinh theo điểm giảm dần, sao cho
tổng chi phí của toàn bộ quá trình nhỏ nhất.

\subsubsection*{Dữ liệu vào}

Tệp \texttt{sorting.in}: dòng đầu là số nguyên $n$ ($2\le n\le 1000$), số
thí sinh. Tiếp theo là $n$ dòng, mỗi dòng chứa một số nguyên không âm $S_i$
($0\le S_i\le 1000000$), điểm của một thí sinh. Có thể giả sử điểm của mọi
người đôi một khác nhau.

\subsubsection*{Dữ liệu ra}

Tệp \texttt{sorting.out}: dòng đầu là số nguyên chỉ số lần chuyển. Mỗi dòng
tiếp theo mô tả một thao tác bằng hai số nguyên $i,j$, nghĩa là chuyển thí
sinh ở vị trí $i$ đến vị trí $j$.

\subsubsection*{Ví dụ}

\begin{verbatim}
sorting.in        sorting.out
5                 2
20                2 1
30                3 5
5
15
10
\end{verbatim}

\subsubsection*{Phân tích}

Để tiện thảo luận, đổi giá trị của mỗi phần tử thành vị trí mục tiêu của nó;
khi đó trạng thái đích là dãy $1,2,\ldots,n$. Giả sử có thêm người thứ $n+1$
đang ở vị trí $n+1$. Thoạt nhìn bài toán không có chỗ để bắt đầu, nhưng phân
tích thêm một chút sẽ thấy: chuyển $i$ đến vị trí $j$ có hai cách nhìn, là
chuyển chủ động hoặc chuyển bị động; trực giác cũng cho thấy nếu một người
phải di chuyển nhiều lần thì không bằng đi thẳng đến nơi.

\noindent\textbf{Giả thuyết 2.1.} Mỗi người nhiều nhất di chuyển một lần. Nếu
di chuyển thì chắc chắn di chuyển đến ngay trước người có số hiệu lớn hơn mình
$1$.

Vì chuyển một người ra sau sẽ làm vị trí của một số người giảm $1$, từ đó giảm
chi phí di chuyển tương lai của họ, nên người có số hiệu lớn hơn nên được xử
lý trước.

\noindent\textbf{Giả thuyết 2.2.} Thực hiện thao tác theo thứ tự số hiệu giảm
dần.

Chứng minh hai giả thuyết trên không phải trọng tâm của bài viết này, nên được
lược bớt ở đây; chứng minh chi tiết nằm trong phụ lục.

\noindent\textbf{Kết luận 2.1.} Thao tác theo thứ tự số hiệu giảm dần. Với mỗi
người $x$ có hai lựa chọn:
\begin{enumerate}
  \item Chuyển đến vị trí ngay trước $x+1$.
  \item Nếu $x$ đang ở trước $x+1$, không di chuyển; về sau sẽ chuyển mọi
  người nằm giữa $x$ và $x+1$ đến trước $x$.
\end{enumerate}

Xét việc chuyển $x$ đang ở vị trí $i$ ra sau đến trước $x+1$ đang ở vị trí
$j$ ($j>i$). Theo Kết luận 2.1, hiện tại các số nhỏ hơn $x+1$ đều chưa di
chuyển; nói cách khác, thứ tự tương đối của mọi số không lớn hơn $x$ giống
ban đầu. Đây là Hệ quả 2.1. Mặt khác, mọi số lớn hơn $x+1$ hoặc đã được
chuyển trước đó, hoặc $x+1$ đã được chuyển đến trước chúng; vì thế mọi số lớn
hơn $x+1$ đều nằm sau $x+1$. Đây là Hệ quả 2.2.

Cần chú ý rằng thao tác chuyển làm thay đổi vị trí, nên $i$ và $j$ không còn
là vị trí ban đầu của $x$ và $x+1$. Ta phải làm gì? Do thứ tự tương đối không
đổi, ta có thể khai thác chính điều này để suy ra vị trí tuyệt đối.

Ví dụ với dãy ban đầu $3,2,5,1,4$, nếu đặt $5$ ra sau $4$ thì $1$ và $4$ đều
tiến lên một vị trí. Nếu thay đổi cách nhìn, cho $5$ và $4$ cùng dùng vị trí
ban đầu $5$, thì $1$ và $4$ xem như không di chuyển, trong khi thứ tự tương
đối của $1,2,3,4$ vẫn giữ nguyên. Thực ra ta có thể dùng thứ tự tương đối để
tính vị trí tuyệt đối: nếu vị trí ban đầu của $x$ là $p_1$, theo Hệ quả 2.1
và 2.2, vị trí thực của $x$ bằng số lượng phần tử nhỏ hơn $x$ trong đoạn ban
đầu $1..p_1$, cộng thêm $1$ cho chính $x$. Nhưng $x+1$ có thể đã di chuyển,
nên vị trí của nó trong dãy ban đầu cần được liệt kê; gọi vị trí đó là $p_2$.

Ký hiệu $c(p,x)$ là vị trí thực của phần tử $x$ sau khi đã xử lý xong các phần
tử $x+1,\ldots,n$, nếu $x$ đang ứng với vị trí $p$ trong dãy ban đầu. Khi đó
$c(p,x)$ bằng số lượng phần tử nhỏ hơn $x$ trong đoạn ban đầu $1..p$, cộng
thêm $1$. Chi phí chuyển $x$ từ vị trí ban đầu $p_1$ đến trước $x+1$ ở vị trí
ban đầu $p_2$ là
\[
  c(p_1,x)+c(p_2,x+1)-1,
\]
vì ta chèn $x$ vào trước $x+1$ nên vị trí đích phải trừ $1$.

Chẳng hạn, trong dãy $3,2,5,1,4$, chuyển $3$ đến trước $4$. Từ hình minh họa
trong bài gốc có thể thấy vị trí thực xuất phát là $i=1$, vị trí đích là
$j=4$. Vị trí thực của $3$ là số phần tử nhỏ hơn $3$ đứng trước $3$ trong dãy
ban đầu, cộng $1$, tức $1$; tương tự, vị trí thực của $4$ là $4$ vì $5$ chắc
chắn ở sau $4$. Chi phí lần chuyển này là $1+4-1$, khớp với tình hình thực tế.
Để không ảnh hưởng đến vị trí của $1$ và $2$, sau khi chuyển, vị trí trong dãy
ban đầu của $3$ trở thành $5$; nghĩa là sau lần chuyển này, các phần tử
$3,4,5$ đều cùng ứng với vị trí ban đầu $5$.

Đặt $f[x][p_2]$ là chi phí nhỏ nhất khi chuyển $x$ đến vị trí $p_2$ trong dãy
ban đầu, và các phần tử $x+1,\ldots,n$ đều đang ở vị trí $p_2$. Gọi
$\mathit{pos}[x]$ là vị trí ban đầu của $x$. Với trường hợp chuyển ra sau, ta
có
\[
  f[x][p_2]=f[x+1][p_2]+c(\mathit{pos}[x],x)+c(p_2,x+1)-1,
  \qquad p_2>\mathit{pos}[x].
\]

Tiếp theo xét trường hợp vị trí đích ở trước vị trí xuất phát. Nếu ban đầu
$4$ và $5$ đều không di chuyển, khi chuyển $3$ đến trước $4$, vị trí của $3$
không phải $c(5,3)$ mà là $c(5,3)+2$. Nếu trước đó quyết định của $5$ là
chuyển ra cuối, còn $4$ không di chuyển, tức $4$ vẫn ở vị trí ban đầu $2$, thì
khi chuyển $3$ đến trước $4$, vị trí của $3$ là $c(5,3)+1$. Nói cách khác,
quyết định trong quá khứ ảnh hưởng đến chi phí di chuyển hiện tại của $3$.
Ta mô phỏng cách xử lý ở bài toán 1: tính ảnh hưởng lên tương lai vào trạng
thái hiện tại. Khi đó trường hợp $j<i$ có thể quy về trường hợp $j>i$.

Ta cần tính ảnh hưởng lên trạng thái tương lai. Nếu $x$ không di chuyển, thì
$x+1$ chắc chắn ở một vị trí nào đó phía sau $x$; gọi vị trí trong dãy ban đầu
của chúng lần lượt là $p_1,p_2$, trong đó $p_1=\mathit{pos}[x]$. Khi quyết
định $5$ không di chuyển, nếu $4$ chuyển đến trước $5$, hoặc $4$ cũng không di
chuyển, thì $3$ đều phải vượt qua cả $4$ và $5$. Nghĩa là trong lần di chuyển
tương lai, vị trí của $3$ là $c(5,3)+(5-3)$. Phần $c(5,3)$ sẽ được tính ở lần
di chuyển tương lai, còn hiện tại chỉ cần cộng thêm phần chắc chắn phát sinh
là $5-3$.

Tóm lại, nếu quyết định $x$ không di chuyển, tuy bước này không sinh chi phí
trực tiếp, ta vẫn phải tính trước một phần chi phí chắc chắn sẽ xuất hiện
trong tương lai: tổng $x-s[k]$ trên các số nhỏ hơn $x$ nằm sau
\(\mathit{pos}[x]\) và trước $p_2$. Viết thành phương trình:
\[
  f[x][\mathit{pos}[x]]
  =
  \min_{p_2>\mathit{pos}[x]}
  \left\{
    f[x+1][p_2]
    +\sum_{\substack{k=\mathit{pos}[x]+1\\ s[k]<x}}^{p_2-1}(x-s[k])
  \right\}.
\]

Tổng hợp lại, phương trình của bài này là
\[
  f[i][j]=f[i+1][j]+c(\mathit{pos}[i],i)+c(j,i+1)-1,
  \qquad j>\mathit{pos}[i],
\]
và
\[
  f[i][\mathit{pos}[i]]
  =
  \min\left\{
    f[i][\mathit{pos}[i]],
    f[i+1][j]
    +\sum_{\substack{k=\mathit{pos}[i]+1\\ s[k]<i}}^{j-1}(i-s[k])
  \right\},
  \qquad j>\mathit{pos}[i].
\]
Đáp án là \(\min_{1\le i\le n} f[1][i]\). Độ phức tạp thời gian là
\(\Oh(n^2)\).

\subsection*{Tiểu kết}

Nhìn lại quá trình giải bài toán 1, trước hết ta phát hiện chi phí của
``hành động'' hiện tại chịu ảnh hưởng của các quyết định quá khứ. Nếu thêm
trạng thái $t$ để biểu diễn ảnh hưởng quá khứ thì số trạng thái quá lớn, nên
ta tính ảnh hưởng đó ngay tại thời điểm quyết định quá khứ và truyền nó qua
trạng thái.

Làm như vậy cần hai điều kiện. Thứ nhất, ảnh hưởng phải là tất yếu. Trong bài
toán 1, chỉ cần ta di chuyển thì quả trứng sẽ rơi. Trong bài toán 2, chỉ cần
ta quyết định không di chuyển, chi phí tương lai sẽ tăng một giá trị cố định.
Nếu bài toán 1 không yêu cầu bắn mọi quả trứng, những quả không bắn trong
tương lai sẽ không chịu ảnh hưởng, và phương pháp trên không còn dùng được.

Thứ hai, ảnh hưởng phải là một hàm của quyết định hiện tại. Ở bài toán 1, hàm
là $-t\sum v$; trong đó $\sum v$ liên quan đến trạng thái và có thể xem như
hằng số, còn $-t$ chính là thời gian di chuyển của quyết định hiện tại. Bài
toán 2 cũng thỏa mãn tính chất này.

Hai điều kiện trên gộp lại có nghĩa là: quyết định hiện tại ảnh hưởng đến chi
phí của ``hành động'' tương lai chỉ thông qua chính quyết định hiện tại.

\section{Ảnh hưởng còn phụ thuộc vào tình huống tương lai}

Một số bài toán không đơn giản như trên. Ảnh hưởng của quyết định hiện tại lên
chi phí ``hành động'' tương lai còn có thể phụ thuộc vào tình huống tương lai.
Khi đó ta thường mở thêm vài chiều trạng thái để giả định quyết định tương
lai, từ đó vẫn tính trước được một phần chi phí như trước; đến khi đi tới
quyết định tương lai thì dùng trực tiếp trạng thái đã giả định.

\subsection{Bài toán 3: Xóa khối (Từ \textit{Nghệ thuật thuật toán và thi tin học})}

\subsubsection*{Đề bài}

Jimmy vừa mua một chiếc máy tính, và giống phần lớn trẻ em, cậu nhanh chóng mê
trò chơi. Gần đây Jimmy đang chơi trò xóa khối. Luật chơi như sau: có $n$ ô
màu xếp thành một hàng. Các khối cùng màu liền kề tạo thành một vùng; nếu hai
khối kề nhau có cùng màu thì chúng thuộc cùng một vùng. Người chơi có thể chọn
tùy ý một vùng để xóa. Nếu vùng đó chứa $x$ khối, người chơi được $x^2$ điểm.
Sau khi một vùng bị xóa, các khối ở bên phải dịch sang trái và nối với phần
bên trái.

\subsubsection*{Phân tích}

Trước hết gộp các khối cùng màu liên tiếp. Dùng cặp $(a,b)$ để biểu thị $b$
khối màu $a$ nối liền nhau.

Mỗi lần xóa một vùng và cộng điểm của vùng đó, nhìn rất giống bài toán gộp đá.
Ta dễ thử đặt $f[i][j]$ là điểm lớn nhất khi xóa hết các vùng từ $i$ đến $j$.
Xét vùng $j$: nếu xóa riêng nó, thì
\[
  f[i][j]=f[i][j-1]+b_j^2.
\]
Còn một khả năng khác: vùng $j$ hợp nhất với một số vùng cùng màu trước đó rồi
mới bị xóa. Nhưng để tính điểm cần biết độ dài sau khi hợp nhất, điều không
thể biết chỉ từ trạng thái hiện tại. Nếu thêm trạng thái để giả định tình
huống quá khứ, ta phải biểu diễn từng khối trong quá khứ đã bị xóa hay chưa,
số trạng thái tăng đến mức không thể chịu được.

Giả sử vẫn muốn dùng cách của hai bài trước: hiện tại xóa một vùng đỏ dài $3$,
trong quá khứ để lại một vùng đỏ dài $3$, tổng điểm khi xóa chung là
$6^2=36$. Nhưng $36$ không thể tách ra để tính rải rác như ở hai bài trên.
Nguyên nhân là hàm điểm là hàm bậc hai; khi quá khứ để lại một vùng đỏ dài
$3$, ảnh hưởng của nó lên việc xóa chung hiện tại còn phụ thuộc vào độ dài
vùng hiện tại.

Vì vậy cần đổi cách nghĩ. Ta muốn hợp nhất vùng $j$ với một số vùng cùng màu
trước nó. Nếu lần hợp nhất này đã được dự liệu từ trước, điểm số đã được tính
trước và truyền đến hiện tại qua trạng thái, thì ta có thể xây dựng phương
trình quy hoạch động. Nói cách khác, với $f[i][j]$ hiện tại ta còn phải giả
định rằng sau một số thao tác xóa, sẽ có $k$ khối cùng màu với $a_j$ nối vào
sau vùng $j$, tạo thành một vùng mới dài $b_j+k$. Đặt
$f[i][j][k]$ là điểm lớn nhất khi xóa/gộp đoạn từ $i$ đến $j$, với giả định
tương lai sẽ có $k$ khối cùng màu $a_j$ nối với $j$.

Vì sao chỉ cần ghi lại tương lai của $j$, mà không cần ghi lại tương lai của
các vùng từ $i$ đến $j-1$?

\noindent\textbf{Chứng minh 3.1.} Nếu vùng $j$ trong tương lai sẽ nối với một
vùng $k_2$, thì các vùng từ $j+1$ đến $k_2-1$ cần bị xóa trước vùng $j$. Vì
thế các vùng trước $j$ chắc chắn không thể nối với các vùng nằm giữa $j$ và
$k_2$. Nếu $j$ trong tương lai nối với $k_1$, còn một điểm $i$ trước $j$ nối
với $k_2$, thì ta có thể tính điểm xóa $i$ trong khi tính trạng thái
$f[i][k_2]$.

Do đó có \textbf{Kết luận 3.1}: với trạng thái $f[i][j][k]$, bất kể thực hiện
thao tác nào, các vùng từ $i$ đến $j-1$ chỉ có thể nối với các vùng trong
đoạn $i..j$, còn vùng $j$ có thể nối với các vùng phía sau $j$.

Từ đó thu được phương trình:
\begin{enumerate}
  \item Nếu xóa trực tiếp vùng thứ $j$ cùng với $k$ khối tương lai nối vào
  sau nó, thì
  \[
    f[i][j][k]=f[i][j-1][0]+(b_j+k)^2.
  \]
  \item Nếu $j$ hợp nhất với một vùng trước đó, giả sử vùng cùng màu với $j$
  là $p$, thì
  \[
    f[i][j][k]=f[i][p][k+b_j]+f[p+1][j-1][0],
    \qquad i\le p<j,\ a_p=a_j.
  \]
\end{enumerate}
Lấy giá trị lớn hơn trong hai loại chuyển trên. Đáp án là $f[1][n][0]$.
Độ phức tạp thời gian là $\Oh(n^4)$.

Ban đầu ta muốn phân bổ chi phí vào từng quyết định ``giữ lại để sau này
xóa'', nhưng phát hiện các chi phí này không độc lập mà liên hệ với nhau:
chi phí hiện tại còn cần biết tình huống tương lai. Vì vậy ta mở thêm một
chiều trạng thái để ghi lại tình huống tương lai.

\subsection{Bài toán 4: Logistics Olympic (NOI 2008, ngày 2)}

\subsubsection*{Đề bài}

Olympic Bắc Kinh 2008 sắp khai mạc, cả nước đang chuẩn bị cho sự kiện trọng
đại này. Để tổ chức Olympic hiệu quả và thành công, việc quy hoạch hệ thống
logistics là không thể thiếu. Hệ thống logistics gồm một số trạm cơ sở, đánh
số từ $1$ đến $N$. Mỗi trạm $i$ có đúng một trạm kế tiếp $S_i$, và có thể có
nhiều trạm tiền nhiệm. Vật tư cần vận chuyển tiếp ở trạm $i$ sẽ được đưa đến
trạm kế tiếp $S_i$. Hiển nhiên trạm kế tiếp của một trạm không thể là chính
nó. Trạm số $1$ gọi là trạm điều khiển; từ bất kỳ trạm nào cũng có thể vận
chuyển vật tư đến trạm điều khiển. Chú ý rằng trạm điều khiển cũng có trạm kế
tiếp, để khi cần có thể lưu thông vật tư.

Trong hệ thống logistics, độ tin cậy cao và chi phí thấp là hai mục tiêu thiết
kế chính. Với trạm $i$, định nghĩa độ tin cậy $R(i)$ như sau. Giả sử trạm
$i$ có $w$ trạm tiền nhiệm $P_1,P_2,\ldots,P_w$, tức các trạm này chọn $i$
làm trạm kế tiếp. Khi đó
\[
  R(i)=C_i+k\sum_{j=1}^{w} R(P_j),
\]
trong đó $C_i$ và $k$ là các hằng số thực dương, và $k<1$. Độ tin cậy của cả
hệ thống tương quan dương với độ tin cậy của trạm điều khiển. Mục tiêu là sửa
đổi hệ thống, tức thay đổi trạm kế tiếp của một số trạm, để $R(1)$ lớn nhất có
thể. Do hạn chế kinh phí, nhiều nhất chỉ được sửa trạm kế tiếp của $m$ trạm,
và không được sửa trạm kế tiếp của trạm điều khiển. Bài toán là: thay đổi
không quá $m$ trạm kế tiếp sao cho $R(1)$ lớn nhất.

\subsubsection*{Dữ liệu vào}

Tệp \texttt{trans.in}: dòng đầu gồm hai số nguyên và một số thực $N,m,k$,
trong đó $N$ là số trạm, $m$ là số trạm kế tiếp được sửa nhiều nhất, còn $k$
là hằng số trong định nghĩa độ tin cậy. Dòng thứ hai gồm $N$ số nguyên
$S_1,S_2,\ldots,S_N$. Dòng thứ ba gồm $N$ số thực dương
$C_1,C_2,\ldots,C_N$.

\subsubsection*{Dữ liệu ra}

Tệp \texttt{trans.out}: in một số thực duy nhất, giá trị lớn nhất của $R(1)$,
chính xác đến hai chữ số sau dấu thập phân.

\subsubsection*{Ví dụ}

\begin{verbatim}
4 1 0.5
2 3 1 3
10.0 10.0 10.0 10.0

30.00
\end{verbatim}

\subsubsection*{Phân tích}

Nối mỗi điểm đến trạm kế tiếp của nó bằng một cạnh. Phân tích đơn giản cho
thấy đồ thị của bài là một số cây có gốc nằm trên một chu trình; các gốc cây
nối với nhau thành chu trình đó.

Với một cây gốc $i$, lặp lại phương trình định nghĩa sẽ thấy $R(i)$ bằng tổng
\[
  \sum_x C_x k^{d(x,i)}
\]
trên mọi điểm $x$ trong cây gốc $i$, trong đó $d(x,i)$ là số cạnh trên đường
từ $x$ đến $i$.

Xét chu trình. Với một chu trình độ dài $3$, chẳng hạn
\[
  R(1)=C_1+kR(2),\qquad
  R(2)=C_2+kR(3),\qquad
  R(3)=C_3+kR(1).
\]
Khử biến được
\[
  R(1)=C_1+kC_2+k^2C_3+k^3R(1),
\]
nên
\[
  R(1)=\frac{C_1+kC_2+k^2C_3}{1-k^3}.
\]
Tức là tổng đóng góp $C_x k^{d(x,1)}$ của mỗi điểm $x$ trên chu trình được
chia cho $1-k^{\mathit{len}}$, trong đó \(\mathit{len}\) là độ dài chu trình.

Với một điểm $x$ nằm trong cây treo vào điểm $i$ trên chu trình, đóng góp của
nó vào $R(1)$ là
\[
  C_x k^{d(x,i)} k^{d(i,1)} = C_x k^{d(x,1)}.
\]
Vì thế
\[
  R(1)=
  \frac{\sum_{i=1}^{n} C_i k^{d(i,1)}}{1-k^{\mathit{len}}}.
\]
Ta liệt kê độ dài chu trình, tức liệt kê cách sửa trạm kế tiếp của một điểm
trên chu trình. Khi đó mẫu số $1-k^{\mathit{len}}$ đã xác định, chỉ cần làm
cho tử số lớn nhất có thể. Vì chu trình đã xác định, ta không được sửa tiếp
trạm kế tiếp của gốc cây; những sửa đổi còn lại chỉ có thể đổi trạm kế tiếp
của một điểm không phải gốc cây. Theo công thức trên, mỗi lần sửa như vậy hiển
nhiên nên đặt trực tiếp trạm kế tiếp của điểm được sửa thành $1$.

Thử dùng quy hoạch động. Với cây gốc $i$, rất tự nhiên đặt trạng thái là giá
trị lớn nhất trong cây gốc $i$ khi đã sửa $j$ lần. Với điểm $i$ có hai lựa
chọn: sửa hoặc không sửa.

Nếu không sửa trạm kế tiếp của $i$, ta phân phối $j$ lần sửa cho các cây con
của $i$, rồi cộng đóng góp của $i$ trong trạng thái hiện tại:
$C_i k^{d(i,1)}$.

Nếu sửa trạm kế tiếp của $i$ thành $1$, đóng góp của cả cây gốc $i$ vào
$R(1)$ phải tính thế nào? Nếu con cháu của $i$ đều chưa từng sửa, đặt trạm kế
tiếp của $i$ thành $1$ sẽ làm khoảng cách từ $i$ và mọi con cháu của $i$ đến
$1$ giảm $2$, tức đóng góp đều chia cho $k^2$. Nhưng nếu trạm kế tiếp của một
điểm con trước đó đã được đặt thành $1$, thì khi ta sửa trạm kế tiếp của $i$
thành $1$, khoảng cách của điểm con đó không hề giảm $2$. Nói cách khác, quyết
định ở điểm con ảnh hưởng đến chi phí khi sửa điểm $i$.

Liệu có thể thêm trạng thái để biểu diễn trong cây $i$ những điểm nào đã sửa,
những điểm nào chưa sửa? Rõ ràng số trạng thái là không thể chấp nhận. Ta cần
tính đóng góp của điểm con vào chi phí sửa điểm $i$ ngay khi quyết định ở điểm
con đó, tức khi quyết định điểm con, cộng luôn chi phí tương lai. Nhưng phải
cộng bao nhiêu? Nếu trạm kế tiếp của $i$ được sửa thành $1$, khoảng cách của
điểm con là $2$. Nếu $i$ không sửa mà một điểm tổ tiên $y$ sửa thành $1$, thì
khoảng cách của điểm con là $3$. Đóng góp của điểm con phụ thuộc vào tổ tiên
được sửa gần nó nhất, vì vậy cần thêm trạng thái biểu diễn tổ tiên được sửa
gần nhất.

Đặt $f[i][j][d]$ là giá trị lớn nhất trong cây gốc $i$, khi sửa $j$ lần, và
khoảng cách từ $i$ đến $1$ là $d$. Xét một con $s$ của $i$: khoảng cách của
$s$ hoặc là $1$ nếu sửa trạm kế tiếp của $s$, hoặc là $d(i)+1$ nếu không sửa.
Đặt
\[
  g[i][j][d]=\max\{f[i][j][d+1],\, f[i][j][1]\}.
\]

Khi đó có các phương trình sau.

\begin{enumerate}
  \item Nếu $i$ không sửa trạm kế tiếp, thì trừ trường hợp bản thân $i$ thật
  sự có khoảng cách $1$ đến $1$, ta không thể có $d=1$. Giả sử các con của
  $i$ là $s_1,s_2,\ldots,s_t$, khi đó
  \[
    f[i][j][d]
    =
    \max_{\substack{j_1+\cdots+j_t=j}}
    \left\{
      \sum_{r=1}^{t} g[s_r][j_r][d]
    \right\}
    +C_i k^d.
  \]
  Phần lấy cực đại có thể thực hiện thêm một lần quy hoạch động kiểu ba lô.
  Gọi \(FF[q][u]\) là giá trị lớn nhất khi đã phân phối $u$ lần sửa cho
  $q$ người con đầu tiên, ta có
  \[
    FF[q][u]=\max_{0\le v\le u}\{FF[q-1][v]+g[s_q][u-v][d]\}.
  \]
  Giá trị cực đại cần lấy là \(FF[t][j]\).

  \item Nếu sửa trạm kế tiếp của $i$, thì
  \[
    f[i][j][1]
    =
    \max_{\substack{j_1+\cdots+j_t=j-1}}
    \left\{
      \sum_{r=1}^{t} g[s_r][j_r][1]
    \right\}
    +C_i k.
  \]
\end{enumerate}

Cuối cùng xét chu trình. Vì chu trình đã được liệt kê và xác định, đáp án cho
phần tử số là
\[
  \max_{\substack{j_1+\cdots+j_{\mathit{len}}=m}}
  \left\{
    f[p_1][j_1][0]+f[p_2][j_2][1]+\cdots+
    f[p_{\mathit{len}}][j_{\mathit{len}}][\mathit{len}-1]
  \right\},
\]
trong đó \(p_1,p_2,\ldots,p_{\mathit{len}}\) là các điểm trên chu trình. Ví
dụ với chu trình độ dài $3$, đáp án là
\[
  \max_{j_1+j_2+j_3=m}
  \{f[1][j_1][0]+f[2][j_2][1]+f[3][j_3][2]\}.
\]
Phần này cũng có thể giải bằng một quy hoạch động nữa. Vì mỗi điểm chỉ bị
tính ba lô một lần, độ phức tạp thời gian là
\(\Oh(\mathit{len}\cdot n^2m^2)\), đủ để vượt qua giới hạn thời gian.

Những năm gần đây xuất hiện nhiều bài tương tự, như IOI 2005 ``Rivers'' và
NOI 2006 ``Network toll''. Các bài này đều quy hoạch động trên cây, đồng thời
chuyển phần chi phí vốn nên tính ở điểm $i$ xuống các hậu duệ của $i$. Khi
quy hoạch các hậu duệ, ta giả định tình huống tương lai và ghi lại bằng trạng
thái; trong bài này, đó là giả định vị trí điểm sửa gần nhất, để khi quy hoạch
đến $i$ có thể trực tiếp dùng quyết định đã giả định trong quá khứ.

\subsection*{Tiểu kết}

Loại bài này có thể xem là phiên bản phức tạp của loại thứ nhất. Ta vẫn phát
hiện quyết định quá khứ ảnh hưởng đến chi phí ``hành động'' hiện tại; ta vẫn
phát hiện nếu ghi lại quyết định quá khứ trong trạng thái hiện tại thì số
trạng thái quá lớn; vì vậy phải tính trước phần ảnh hưởng đó trong quá khứ.
Tuy nhiên ảnh hưởng của quyết định hiện tại lên chi phí ``hành động'' tương
lai còn liên quan đến tình huống tương lai. Khi đó ta giả định tình huống
tương lai, rồi truyền các ảnh hưởng khác nhau đến tương lai theo các trạng
thái khác nhau.

\section{Tổng kết}

Bốn ví dụ trong bài đều có phương trình không dễ xây dựng. Nguyên nhân là chi
phí của ``hành động'' hiện tại chịu ảnh hưởng của các quyết định quá khứ; còn
nếu thêm trạng thái để giả định quá khứ ở hiện tại thì số trạng thái quá lớn.
Vì vậy ta thay đổi ``quan niệm thời gian'': thay vì nhìn hiện tại từ quá khứ,
ta nhìn tương lai từ hiện tại, và tính ảnh hưởng của quyết định hiện tại lên
tương lai vào trạng thái hiện tại.

Với loại bài thứ nhất, chỉ cần xem ảnh hưởng lên tương lai là chi phí của
quyết định hiện tại và lưu trong trạng thái hiện tại. Với loại bài thứ hai,
cần giả định tình huống tương lai, lưu ảnh hưởng của các tình huống tương lai
khác nhau trong các trạng thái khác nhau. Có thể hiểu là truyền ảnh hưởng đến
tương lai theo những con đường khác nhau, để khi đến quyết định tương lai thì
dùng trực tiếp.

Dĩ nhiên hai loại bài này có biến hóa vô tận. Nhưng chỉ cần giỏi phát hiện và
dám đổi mới, ta sẽ xây dựng được phương trình.

\section*{Lời cảm ơn}

Xin chân thành cảm ơn các thầy Cao Lợi Quốc và Chu Tổ Tùng đã hướng dẫn, giúp
đỡ tôi khi viết bài này. Xin chân thành cảm ơn các bạn Tầm Vân Ba, Lý Viễn
Thao, Ngô Không, Dịch Thiến, Đặng Phương Đinh, Ngô Nhất Phàm và nhiều bạn khác
đã giúp đỡ tôi.

\section*{Tài liệu tham khảo}

\begin{enumerate}[label={[\Alph*]}]
  \item \textit{Nghệ thuật thuật toán và thi tin học}, Lưu Nhữ Giai và Hoàng
  Lượng, Nhà xuất bản Đại học Thanh Hoa.
  \item Bài giảng WC 2008 của Tôn Huy.
  \item Lời giải chính thức BOI 2007.
\end{enumerate}

\appendix

\section{Chứng minh cho Giả thuyết 2.1 và 2.2}

Phần chứng minh sau lấy từ lời giải chính thức BOI 2007.

Xét thí sinh \(p_{\min}\) có điểm nhỏ nhất; cuối cùng người này phải đứng ở vị
trí \(n\).

\noindent\textbf{Mệnh đề 1.} Nếu \(p_{\min}\) được di chuyển, phương án tối ưu
là chuyển thẳng người này đến cuối danh sách.

\noindent\textit{Chứng minh.} Hiển nhiên chuyển \(p_{\min}\) về phía trước chỉ
làm tăng chi phí của những người khác cần di chuyển. Giả sử người này được
chuyển ra sau đến một vị trí \(i<n\), tức tiết kiệm \(n-i\) bước so với chuyển
đến cuối. Khi đó có hai khả năng: hoặc \(n-i\) người phía sau người này trong
danh sách phải được di chuyển, và chi phí của họ đều tăng thêm \(1\) so với
trường hợp \(p_{\min}\) đã được chuyển đến cuối; hoặc \(p_{\min}\) phải di
chuyển thêm lần nữa, điều hiển nhiên không tối ưu.

\noindent\textbf{Mệnh đề 2.} Nếu \(p_{\min}\) được di chuyển, ta có thể thực
hiện việc này ở bước đầu tiên.

\noindent\textit{Chứng minh.} Từ Mệnh đề 1, ta đã biết phần thứ hai của chi phí
di chuyển là cố định. Vì vậy lý do duy nhất để không chuyển \(p_{\min}\) trước
tiên là vị trí hiện tại của người này có thể giảm do các thí sinh khác được
chuyển đi. Nhưng chi phí di chuyển của những thí sinh đó lại tăng thêm \(1\)
vì \(p_{\min}\) chưa được chuyển đến cuối.

Nếu \(p_{\min}\) được di chuyển, bài toán giảm thành cùng một bài toán với
\(n-1\) thí sinh. Nếu \(p_{\min}\) không di chuyển, ta có một bài toán tương
tự với \(n-1\) thí sinh, kèm ràng buộc bổ sung rằng mọi thí sinh đứng sau
\(p_{\min}\) phải được di chuyển trước \(p_{\min}\).

Gọi \(p'_{\min}\) là thí sinh có điểm nhỏ nhất trong \(n-1\) thí sinh còn lại.
Ta thấy không có ích gì khi chuyển \(p'_{\min}\) ra sau \(p_{\min}\), vì như
vậy sẽ phải chuyển người này thêm lần nữa, và việc vượt qua các thí sinh vẫn
chưa di chuyển cũng không đem lại lợi ích: với mỗi thí sinh như vậy, ta giảm
chi phí di chuyển của họ đi \(1\), nhưng lại tăng chi phí di chuyển của
\(p'_{\min}\) ít nhất \(1\). Nếu \(p'_{\min}\) hiện đứng trước \(p_{\min}\)
trong danh sách, bằng lập luận tương tự Mệnh đề 1, nếu \(p'_{\min}\) được di
chuyển thì tối ưu là chuyển thẳng đến trước \(p_{\min}\). Mệnh đề 2 vẫn còn
đúng.

Trường hợp còn lại là \(p'_{\min}\) hiện đứng sau \(p_{\min}\) và phải được
chuyển về phía trước. Nếu có các thí sinh nằm giữa \(p_{\min}\) và
\(p'_{\min}\), ta có thể cân nhắc chuyển họ trước. Nếu chuyển \(p'_{\min}\)
trước, vị trí xuất phát của các thí sinh ở giữa tăng thêm \(1\). Nhưng nếu
chuyển một thí sinh ở giữa trước, vị trí đích mà \(p'_{\min}\) cần chuyển đến
cũng tăng thêm \(1\). Do đó chuyển \(p'_{\min}\) trước không thể tệ hơn.

Vì vậy, bằng quy nạp và các lập luận trên, ta chứng minh được rằng mỗi thí
sinh được di chuyển nhiều nhất một lần, và các thí sinh có thể được di chuyển
theo thứ tự điểm tăng dần, tức theo số hiệu mục tiêu giảm dần trong cách đánh
số của phần phân tích.

\section{Mã nguồn}

Phụ lục B của bản PDF gốc chỉ hiện các tiêu đề ``mã nguồn bài toán 1'', ``mã
nguồn bài toán 2'' và ``mã nguồn bài toán 4'' trong kết quả trích xuất văn
bản; nội dung chương trình không được \texttt{pdftotext} trích xuất thành văn
bản có thể dịch.

\end{document}
